\section{Method}\label{sec:method}

A rig $\mathcal{R}=(\mathcal{T},\mathbf{p}^{\mathrm{rest}},\mathcal{D})$ is a kinematic tree
$\mathcal{T}$ over $J$ joints, a rest pose $\mathbf{p}^{\mathrm{rest}}\in\mathbb{R}^{J\times3}$ and
a natural-language description $\mathcal{D}=\{d_j\}_{j=1}^{J}$ of its joints. A motion is a tensor
$\mathbf{x}\in\mathbb{R}^{L\times J\times 17}$ in the \rep{} representation of
\S\ref{sec:method:rep}. Given $\mathcal{R}$ and a caption $c$, \method{} samples
$\mathbf{x}\sim p_\theta(\mathbf{x}\mid\mathcal{R},c)$ with one transformer shared by every rig
(\S\ref{sec:method:model}), trained by a calibrated grouped flow-matching objective
(\S\ref{sec:method:objective}) and sampled in 20 Euler steps
(\S\ref{sec:method:sampling}). A rig outside the training library is served by the same
forward pass and, when a few clips are available, by a per-rig LoRA
(\S\ref{sec:method:lora}); all retrieval metrics come from a skeleton-aware evaluator
(\S\ref{sec:method:eval}).

\subsection{\rep{}: a per-joint representation for any rig}\label{sec:method:rep}

\paragraph{Channels.} Every joint carries the same 17 channels (Table~\ref{tab:channels} in App.~\ref{app:rep}). Positions
are expressed relative to the root's own ground-plane track, with height absolute, so
that the body's pose and its travel are separated the way animators separate them. Rotations are
continuous 6D~\citep{zhou2019continuity} \emph{deltas} from the joint's rest orientation in the global
frame: ``no motion'' is the identity rotation, $[1,0,0,0,1,0]$ before normalisation, on every rig,
so that the rest frame means the same thing on every rig when it serves as the skeleton prompt
(\S\ref{sec:method:model}). Velocities are supervision only and are never
integrated at decode time, and the root comes from the trajectory channels. Every frame therefore decodes on its
own, without accumulated drift, by either of two routes that do not read each other: the position channels, or
forward kinematics from the rotation channels. Their disagreement is the FK--pose gap this paper reports --- a property
the representation exposes, not a metric imposed on it. Everything is expressed in one canonical frame ($+Y$ up, rest forward $+Z$) obtained
by a rigid per-rig re-alignment to the foot support plane, after which forward kinematics is
unchanged to numerical precision (App.~\ref{app:rep}). Building \rep{} for a new rig needs only its
kinematic tree, rest pose and animation: there is no template skeleton, no retargeting and no
inverse-kinematics fitting, and the same channel definition converts every
Truebones rig used in \S\ref{sec:method:lora}.

\paragraph{Per-cell normalisation.} Rigs differ in size, proportion and convention by orders of
magnitude, so every (rig, joint, channel) cell is standardised with its own mean and standard
deviation, the latter floored at $\epsilon_\sigma$ (\S\ref{sec:exp:setup}), and cells that are constant across
a rig's library, such as never-firing contact flags, are removed from supervision (App.~\ref{app:rep}). The
statistics are a property of the rig, computed from its motion library and supplied with it at deployment,
just as motion-transfer tools require example motions of the target
character~\citep{chen2025motion2motion}. The library, its split and the conversion are described in
\S\ref{sec:exp:setup} and App.~\ref{app:rep}.

\subsection{In-context skeleton conditioning}\label{sec:method:model}

\paragraph{Tokens and layout.} One token per frame and joint (Figure~\ref{fig:teaser}, \emph{any rig
$\to$ tokens}): the 17 channels of a cell are projected to width $D$ and receive a frame-index and a
joint-slot embedding. The sequence is $[\,\mathbf{x}^{\mathrm{demo}}\mid\mathbf{x}_t^{\mathrm{tgt}}\,]$
along time: one clean frame holding the rig's analytic rest pose, never noised, followed by the $L$
target frames (\S\ref{sec:exp:setup}), which carry the flow interpolant and a learned marker and are the
only frames denoised and scored. Under the rest-delta parameterisation the rest frame's rotations are the identity
on every rig, so it conveys proportions, joint count and graph rather than motion; after per-cell
normalisation it also carries the rig's statistics.

\paragraph{Three signals identify the skeleton.} Nothing in the network is indexed by rig; what it knows
about a body arrives as input (Figure~\ref{fig:teaser}, \emph{token embedding}). \emph{Geometry}: the
rest frame above. \emph{Identity}: each joint description $d_j$ is embedded once with frozen
LLM2Vec~\citep{behnamghader2024llm2vec} (4096-d), projected and added to every token of joint $j$.
We use descriptions rather than bone names because ``front left upper leg'' is the same concept on a horse
and a lizard whereas \texttt{L\_Femur\_02} and \texttt{b\_hindleg\_l\_1} are not, and a new rig's
naming scheme need never have been seen. \emph{Place in the tree}: an 8-d structural vector per joint
(rest-offset direction, bone length, tree depth, root-path length, child count, leaf flag) passes
through a zero-initialised MLP and is added to the same tokens.

\paragraph{The graph enters as an attention bias.} Each of the $N_b$ blocks factorises attention over the
two token axes (Figure~\ref{fig:teaser}, \emph{transformer}): temporal attention, one joint over its
frames; spatial attention, one frame over its joints; then a $4\times$ MLP. Only the spatial attention
sees the kinematic tree, as an additive per-head bias on the logits between joints $i$ and $j$,
$b^{(h)}_{ij}=-g_{ij}+u^{(h)}_{a_{ij}}+d^{(h)}_{s_{ij}}$, where $g_{ij}$ is the hop distance on the tree,
clipped at $g_{\max}$, $a_{ij},s_{ij}$ are the hops up and down to the lowest common ancestor, clipped at
$a_{\max}$ (\S\ref{sec:exp:setup}), and $u,d$ are zero-initialised per-head tables. Nearer joints attend
more by default, and the tables let a head tell parent from child, and a joint on the same limb from one
across the body (App.~\ref{app:model}). The
tree is a soft prior on a fully connected joint attention, not a fixed layer, so one network serves
any joint count and topology.

\paragraph{Text and time modulate.} The conditioning vector
$\mathbf{c}=\mathrm{MLP}(t)+\mathrm{MLP}(\mathbf{e}_c)$, with $\mathbf{e}_c$ the pooled LLM2Vec caption
embedding, drives the AdaLN~\citep{peebles2023dit} shift, scale and gate of every sub-layer (shift and scale
only on the output head), so text reaches the model only through modulation and the token stream stays
kinematic. The caption is dropped with probability $p_{\mathrm{drop}}$ (\S\ref{sec:exp:setup}) for
classifier-free guidance; the rest-pose prompt is never dropped.

\paragraph{Stability.} Queries and keys are RMS-normalised per head before the dot
product~\citep{dehghani2023vit22b}; without it, every attempt at 0.3B diverged (App.~\ref{app:ablation_notes}).

\subsection{Calibrated grouped flow-matching objective}\label{sec:method:objective}

\paragraph{Flow matching with $\mathbf{x}$-prediction.} We train on the straight path
$\mathbf{x}_t=(1-t)\,\mathbf{x}_0+t\,\mathbf{x}_1$ with $\mathbf{x}_0\sim\mathcal{N}(0,I)$ and
$t\sim\mathcal{U}(0,1)$~\citep{lipman2023flow,liu2023rectified}; the network predicts the clean
sample $\hat{\mathbf{x}}_1$, and the velocity-space weight of \citet{li2025jit} is applied,
\begin{equation}
  \ell(\mathbf{x}_t,t)=\frac{w(t)}{\mathbb{E}[w]}\;\mathrm{Huber}_{\delta}\!\left(\hat{\mathbf{x}}_1-\mathbf{x}_1\right),
  \qquad w(t)=\frac{1}{\max(1-t,\,\sigma_{\min})^{2}},
  \label{eq:elem}
\end{equation}
with $\mathrm{Huber}_{\delta}(r)=r^{2}$ for $|r|\le\delta$ and $\delta(2|r|-\delta)$ beyond, applied per cell;
$\delta$ sits below the 99.99th percentile of the normalised targets (App.~\ref{app:calib}), $\sigma_{\min}$
and $\delta$ are given in \S\ref{sec:exp:setup}, and $t$ is sampled uniformly (App.~\ref{app:ablation_notes}).

\paragraph{Grouped loss.} A rig with 102 joints has three times the body cells of one with 34 and
the same single root row, so a plain mean over cells would drown the root trajectory, heading and
contact channels---the ones that decide whether a walk is a walk---in body rotations on large rigs.
We partition the channels into nine groups $g$: position, rotation and velocity, each for the root and
for the body, plus contact, root track and heading, and use
\begin{equation}
  \mathcal{L}=\sum_{g}\gamma_g\sqrt{\frac{N_g}{N}}\;\bar{\ell}_g,
  \label{eq:grouped}
\end{equation}
where $\bar{\ell}_g$ is the masked mean of Eq.~(\ref{eq:elem}) over the $N_g$ valid cells of group
$g$ in the batch and $N=\sum_g N_g$; the $\gamma_g$ are set by the calibration below, following the
target-share weighting of Kimodo~\citep{rempe2026kimodo}.
Under this weighting the explicit $N_g$ dependence of the output-gradient energy of group $g$
cancels: it is proportional to $\gamma_g^{2}Q_g$, where $Q_g$ is the group's mean weighted squared
residual, so a group's weight in $\partial\mathcal{L}/\partial\hat{\mathbf{x}}_1$ no longer grows
with the number of cells the group spans in the batch (App.~\ref{app:calib}).

\paragraph{Calibration.} The $\gamma_g$ are not hand-tuned. A target share profile over channel
families is fixed in advance, the group energies $E_g$ are measured over one pass of the training set
under exactly the training crops and masks, and the weights are solved as
$\gamma=\sqrt{\text{share}/E}$, with the rotation family anchored at $\gamma_{\mathrm{rot}}$
(\S\ref{sec:exp:setup}) and the others solved relative to it. A \emph{mechanism check} then takes a few
optimiser-free steps of a small model of the same recipe and verifies that every family's measured share of
$\|\partial\mathcal{L}/\partial\hat{\mathbf{x}}_1\|^{2}$ matches the $\gamma_g^{2}Q_g$ prediction
(App.~\ref{app:calib}). It certifies the wiring of the weights, not that the profile fixed in advance is
realised during training: the realised share moves with the error energies, which evolve.

\paragraph{Auxiliary terms.} Four terms are added with weights $\lambda_{\mathrm{FK}}$,
$\lambda_{\mathrm{vel}}$, $\lambda_{\mathrm{lock}}$ and $\lambda_{\mathrm{acc}}$ (\S\ref{sec:exp:setup}), none
carrying the time weight $w(t)$: a forward-kinematics consistency term between predicted positions and
positions forward-kinematised from predicted rotations, a velocity term on world joint displacements, a
foot-lock term gated by the contact channel, and an acceleration-matching term (App.~\ref{app:calib}).

\subsection{Sampling}\label{sec:method:sampling}
Samples follow the straight path with Euler steps
$\mathbf{x}\leftarrow\mathbf{x}+(\hat{\mathbf{x}}_1-\mathbf{x})/(S-i)$ over $S$ steps; invalid cells are
projected back to zero after every step, and guidance is applied to the caption only,
$\hat{\mathbf{x}}_1=\hat{\mathbf{x}}_1^{\varnothing}+s\,(\hat{\mathbf{x}}_1^{c}-\hat{\mathbf{x}}_1^{\varnothing})$
with guidance scale $s$ (\S\ref{sec:exp:setup}). Twenty steps is the operating point used throughout; ten give the same retrieval at half the
cost, and App.~\ref{app:cost} gives the latency and quality of both.

\subsection{Extending the model to a rig it has never seen}\label{sec:method:lora}
\paragraph{Zero-training arm.} An unseen rig is served by the frozen backbone as soon as its rest
pose, joint descriptions and per-cell statistics (from its own clips) are supplied; on the Truebones
rigs this yields motion with far more jitter than the real clips, since wing chains and eight-legged
gaits are absent from the training prior.

\paragraph{Per-rig LoRA.} When roughly ten clips of the new rig exist, we attach rank-64 or rank-128
adapters~\citep{hu2022lora} (scale 1) to the attention projections, the MLPs and the conditioning
paths (joint-semantics projection, caption MLP, structural MLP, AdaLN), leaving the input and output
projections and the timestep MLP frozen; the adapters hold 26.8M (r64) or 53.6M (r128) parameters
and are folded back into the weights at export, so every consumer loads a plain checkpoint. The
new rig's data view is built by the corpus script (joint pruning to the deformation-bone convention,
corpus-style captions, an 80/20 split) and the adapter is trained with the backbone
recipe for 300--900 epochs (App.~\ref{app:lora}); a single adapter shared by a group of rigs
removes the same jitter excess on all of them, and matches a per-rig adapter on the one rig where both were
trained (App.~\ref{app:extend_more}).

\subsection{Skeleton-aware evaluation}\label{sec:method:eval}
\paragraph{Evaluator.} Retrieval metrics require a text--motion embedding, and existing ones are
tied to one skeleton. We train two towers into a shared 512-d space: a frozen DistilBERT text tower
with a trainable pooling head, and a graph-aware motion tower (6 graph layers, 4 temporal layers)
over the \rep{} channels without contact, with a symmetric InfoNCE objective. The
evaluator sees only real data, never generated samples, and shares no weights with the generator.

\paragraph{Protocol.} Text$\to$motion R@1/2/3 in pools of 64, matching score and FID in the
evaluator space, each next to the evaluator's score on the real clips of the same pools (with shuffled
captions the evaluator scores R@1 \num{0.035} in pools of 32, chance \num{0.031}); generation is fixed at 20 steps,
$s{=}2$ at seed 42 and sharded across GPUs under a hashed generation plan with provenance checks.
Rigs outside the library, for which no evaluator exists, are scored geometrically
(\S\ref{sec:exp:setup}).
